\documentclass[acmsmall,screen,nonacm]{acmart}
\usepackage{lambdapcc}

\title{Capture Checking for Path-Dependent Pairs}
\author{Oliver Bra\v{c}evac}
\author{Martin Odersky}
\renewcommand{\shortauthors}{Bra\v{c}evac and Odersky}
\authorsaddresses{}

\begin{document}

\begin{abstract}
We extend a small calculus of path-dependent types based on immutable
dependent pairs with capturing types, use sets, and abstract capture-set
members. Capture sets contain paths of arbitrary length. A pair member may
describe a term, a type interval, or a capture-set interval, and its signature
may depend on the pair's first component. Subcapturing places the singleton
capture set of a term path below the capture set of its synthesized type, relates an
abstract capture-set selection to its declared bounds, and contracts
projections and term-member selections to their root.

The soundness proof interprets subtyping derivations as store-indexed
coercions. An annotated store records, for each allocated value, the capture
set assigned by its introduction rule. Value typing retains this set
and a subcapturing derivation to the capture set of the assigned type.
Coercion action preserves the assigned capture set through dependent functions,
pairs, and abstract members. The general pair rule is handled by applying its first-component
coercion at the stored component and instantiating its member coercion at the
same location. Coercion action is well founded in the allocation order of
the immutable store.

The mechanization proves progress, preservation, and non-stuckness for
closed programs.  It also proves that both operands inspected by every
application are covered by the source use set, transported through preceding
allocations. For a returned value, the capture set assigned by its
introduction rule subcaptures the capture set of the result type.
\end{abstract}

\keywords{path-dependent types, capture checking, dependent pairs,
subtyping, semantic coercions, type safety}

\maketitle

\section{Introduction}

Path-dependent types allow the type of a program phrase to mention a stable
path. Capture checking attaches a finite description of capabilities to a
value and tracks the capabilities used while evaluating a term. Both
mechanisms become more useful when paths may be longer than variables: a
component selected from a structure can carry its own type and its own
capture information.

Capabilities are named by stable term paths in this calculus. Capture sets
over-approximate the capability paths retained by values, while use sets
over-approximate those used during evaluation.

The calculus studied here, \lambdapcc{}, combines these ideas in a small
call-by-value language of immutable dependent pairs. It retains the path,
singleton, abstract-type, function, and general pair-subtyping mechanisms of
the underlying path calculus. It adds three ingredients from capture
checking. First, every type is a capturing type.
Second, term typing produces a \emph{use set}, an upper bound on the
capabilities used to evaluate the term. Third, pairs may contain abstract
capture-set members with lower and upper bounds, in parallel with abstract
type members.

A capturing type \(S\capt C\) consists of a shape type \(S\) and a capture
set \(C\), which upper-bounds the capabilities that a value may capture. A typing
judgment
\[
  \typ{C}{\Gamma}{t}{T}
\]
assigns result type \(T\) and use set \(C\) to \(t\). Values have an empty
use set because their evaluation is complete. For an abstraction, the use set
of the body becomes the capture set of the function type. This separation gives the usual distinction
between capabilities used now and capabilities retained for later use
\citep{xu2025box,boruch2023capture}.

Capture sets may contain a capability path \(p\) and an abstract capture-set
selection \(p.a\). We write the former as the singleton \(\set p\); the
latter is itself a capture-set expression. The two forms have different
roles. If \(p\) synthesizes \(S\capt C\), subcapturing derives
\(\set p<:C\). If the capture-set member \(p.a\) has bounds \(L..U\), then
\(L<:p.a<:U\). Projections and term-member selections contract to their
owner, for example \(\set{p.a}<:\set p\). There is no corresponding rule
\(p.a<:\set p\): a capture-set member is a static definition whose contents
need not be owned by its receiver. This distinction prevents an arbitrary
concrete capture-set definition from being collapsed through the receiver's
capture set.

A member signature \(\mu\) is a capturing type, a type interval, or a
capture-set interval. The main subtype of interest is the general pair rule:
\begingroup
\rulesetup
\infrule[\ruleuse{ss-pair}]
  {\Gamma\vdash T<:T'
   \andalso
   \Gamma,x:T\vdash\mu<:\mu'}
  {\Gamma\vdash
    \pairshape{x}{T}{a}{\mu}<:
    \pairshape{x}{T'}{a}{\mu'}}
\endgroup
The premise under \(x:T\) cannot be compiled to a closed semantic
coercion until a concrete pair supplies the location denoted by \(x\).
The proof retains this premise as a suspended member premise. When a pair at location
\(z\) is observed, its runtime value exposes a first component \(y\) and a
member referent. The proof coerces \(y\), instantiates the suspended premise with
\(x\mapsto y\), and applies the resulting coercion to that referent.

The annotated store records each allocated value together with the capture
set \(Q\) assigned by its introduction rule. Runtime typing of that value at
\(S\capt C\) retains a derivation \(Q<:C\). For an abstraction, \(Q\) is
obtained from the body's use set; for a pair, it is determined by the stored
components. The premise of a value-introduction rule may already contain
subsumption, so \(Q\) need not be least. Runtime term typing also retains the
use set of the current term. Retaining both sets supports separate results for type
safety, application coverage, and the capture-set bound on returned values.

The rest of the paper presents the complete calculus, its operational
semantics, the coercion interpretation, and the soundness argument. The
presentation follows the formal judgments rather than the auxiliary
organization of the mechanization.

\section{The calculus}
\label{sec:calculus}

\subsection{Notation and syntactic categories}

Variables are written \(x,y,z\), and labels are written \(a,b\). Labels
inhabit one syntactic namespace. Path synthesis determines whether a
selection denotes a term, type, or capture-set member. Paths \(p,q\) consist
of variables, first projections, and labelled selections. A \emph{term path}
is a path that synthesizes a capturing type; only term paths may occur as
terms or singleton capture sets.

Capture sets are written \(C,D\), shape types \(R,S\), types \(T,U\),
and member signatures \(\mu,\nu\). A member signature is a
capturing type, a shape-type interval, or a capture-set interval. We omit an
empty capture set, writing \(S\) for \(S\capt\set{}\). The singleton
shape type of a path is \(\sing p\), whereas \(\set p\) is a singleton capture
set. The binder in a function scopes over its codomain. The binder in
\(\pairshape{x}{T}{a}{\mu}\) scopes over \(\mu\) and denotes the first
component. Capture-avoiding opening is written \(\subst{X}{p}{x}\) for any
type, shape type, signature, or capture set \(X\).

Objects well scoped in \(\Gamma\) may also be used under an extension
\(\Gamma,x:T\); this weakening is implicit. Primes and subscripts denote
further metavariables of the same category. In rules that apply uniformly to
type and capture-set intervals, \(L\) and \(U\) denote the lower and upper
bound; the interval form determines whether they range over shape types or
capture sets.

Terms are in monadic normal form: applications take paths, and pairs store a
first component together with a term, shape-type, or capture-set definition.
Shape-type and capture-set definitions are retained only as static lookup
metadata. Their paths are not captured capabilities, so pair introduction
records only term components in the capture set.

The static judgments are path synthesis \(\Gamma\vdash p\Rightarrow\mu\),
well-formedness \(\Gamma\vdash X\ \wf\), subcapturing
\(\Gamma\vdash C<:D\), subtyping \(\Gamma\vdash T<:U\),
\(\Gamma\vdash S<:R\), and \(\Gamma\vdash\mu<:\nu\), and term typing
\(\typ{C}{\Gamma}{t}{T}\).

\clearpage
\begingroup
\setlength{\abovedisplayskip}{5pt plus 1pt minus 1pt}
\setlength{\belowdisplayskip}{5pt plus 1pt minus 1pt}
\setlength{\intextsep}{5pt}
\captionsetup{skip=6pt}
\rulesetup
\setlength{\afterruleskip}{0.65em plus 0.1em minus 0.05em}
\renewcommand{\plateheader}[2]{%
  \par\addvspace{0.45em}%
  \noindent\textbf{#1}%
  \if\relax\detokenize{#2}\relax\else
    \hfill{\setlength{\fboxsep}{3pt}\setlength{\fboxrule}{0.5pt}%
      \fbox{\(\strut #2\)}}%
  \fi\par
  \nobreak\addvspace{0.45em}%
}

\begin{figure}[H]
\plateheader{Source syntax}{}
\noindent
\begingroup
\renewcommand{\arraystretch}{0.94}
\begin{minipage}[t]{0.48\linewidth}
\begin{tabular*}{\linewidth}{@{}r@{\ }l@{\extracolsep{\fill}}r@{}}
\(p,q ::= \) && \textbf{Path}\\
 & \(x\) & variable\\
 & \(p.1\) & first projection\\
 & \(p.a\) & member selection\\[0.35em]
\(C,D ::= \) && \textbf{Capture set}\\
 & \(\set{}\) & empty\\
 & \(C\cup D\) & union\\
 & \(\set p\) & capability path\\
 & \(p.a\) & capture-set selection\\[0.35em]
\(T,U ::= \) && \textbf{Type}\\
 & \(S\capt C\) & capturing\\[0.35em]
\(R,S ::= \) && \textbf{Shape type}\\
 & \(\top\mid\bot\) & extrema\\
 & \((x:T)\to U\) & function\\
 & \(\pairshape{x}{T}{a}{\mu}\) & pair\\
 & \(\sing p\) & singleton\\
 & \(p.a\) & type selection
\end{tabular*}
\end{minipage}\hfill
\begin{minipage}[t]{0.48\linewidth}
\begin{tabular*}{\linewidth}{@{}r@{\ }l@{\extracolsep{\fill}}r@{}}
\(\mu,\nu ::= \) && \textbf{Member signature}\\
 & \(T\) & term member\\
 & \(S..R\) & type interval\\
 & \(C..D\) & capture-set interval\\[0.35em]
\(\delta ::= \) && \textbf{Member definition}\\
 & \(z\) & term definition\\
 & \(S\) & type definition\\
 & \(C\) & capture-set definition\\[0.35em]
\(t,u ::= \) && \textbf{Term}\\
 & \(p\) & path\\
 & \(\lambda(x:T).t\) & abstraction\\
 & \(\pairval{y}{a}{\delta}\) & pair\\
 & \(p\,q\) & application\\
 & \(\LET x=t\ \IN u\) & let\\[0.35em]
\(v ::= \) && \textbf{Value}\\
 & \(\lambda(x:T).t\mid\pairval{y}{a}{\delta}\) &\\[0.35em]
\(\Gamma ::= \) && \textbf{Context}\\
 & \(\ctxempty\mid\Gamma,x:T\) &
\end{tabular*}
\end{minipage}
\endgroup

\plateheader{Path synthesis}{\Gamma\vdash p\Rightarrow\mu}
\noindent
\begin{minipage}[t]{0.47\linewidth}
\infrule[\ruledef{p-var}]
  {\Gamma(x)=T}
  {\Gamma\vdash x\Rightarrow T}

\infrule[\ruledef{p-fst}]
  {\Gamma\vdash p\Rightarrow
    \pairshape{x}{T}{a}{\mu}\capt C}
  {\Gamma\vdash p.1\Rightarrow T}
\end{minipage}\hfill
\begin{minipage}[t]{0.51\linewidth}
\infrule[\ruledef{p-sel}]
  {\Gamma\vdash p\Rightarrow
    \pairshape{x}{T}{a}{\mu}\capt C}
  {\Gamma\vdash p.a\Rightarrow\subst{\mu}{p.1}{x}}

\infrule[\ruledef{p-miss}]
  {\Gamma\vdash p\Rightarrow
     \pairshape{x}{T}{b}{\nu}\capt C
   \\
   \Gamma\vdash p.1.a\Rightarrow\mu
   \andalso a\ne b}
  {\Gamma\vdash p.a\Rightarrow\mu}
\end{minipage}

\plateheader{Capture-set and member-signature well-formedness}
 {\Gamma\vdash C\ \wf\qquad\Gamma\vdash\mu\ \wf}
\noindent
\begin{minipage}[t]{0.48\linewidth}
\noindent\textit{Capture sets}\par\nobreak
\begin{minipage}[t]{0.46\linewidth}
\infax[\ruledef{wfc-empty}]
  {\Gamma\vdash\set{}\ \wf}

\infrule[\ruledef{wfc-union}]
  {\Gamma\vdash C\ \wf\quad\Gamma\vdash D\ \wf}
  {\Gamma\vdash C\cup D\ \wf}
\end{minipage}\hfill
\begin{minipage}[t]{0.51\linewidth}
\infrule[\ruledef{wfc-path}]
  {\Gamma\vdash p\Rightarrow T}
  {\Gamma\vdash\set p\ \wf}

\infrule[\ruledef{wfc-select}]
  {\Gamma\vdash p.a\Rightarrow C..D
   \\
   \Gamma\vdash C<:D}
  {\Gamma\vdash p.a\ \wf}
\end{minipage}
\end{minipage}\hfill
\begin{minipage}[t]{0.50\linewidth}
\noindent\textit{Member signatures}\par\nobreak
\begin{minipage}[t]{0.31\linewidth}
\infrule[\ruledef{wfm-term}]
  {\Gamma\vdash T\ \wf}
  {\Gamma\vdash T\ \wf}
\end{minipage}\hfill
\begin{minipage}[t]{0.66\linewidth}
\infrule[\ruledef{wfm-type}]
  {\Gamma\vdash S\ \wf\quad\Gamma\vdash R\ \wf
   \quad\Gamma\vdash S<:R}
  {\Gamma\vdash S..R\ \wf}

\infrule[\ruledef{wfm-capture}]
  {\Gamma\vdash C\ \wf\quad\Gamma\vdash D\ \wf
   \quad\Gamma\vdash C<:D}
  {\Gamma\vdash C..D\ \wf}
\end{minipage}
\end{minipage}

\plateheader{Type and shape-type well-formedness}
  {\Gamma\vdash T\ \wf\quad\Gamma\vdash S\ \wf}
\noindent
\begin{minipage}[t]{0.31\linewidth}
\infrule[\ruledef{wft-capt}]
  {\Gamma\vdash C\ \wf\andalso\Gamma\vdash S\ \wf}
  {\Gamma\vdash S\capt C\ \wf}

\infax[\ruledef{wfs-bot}]
  {\Gamma\vdash\bot\ \wf}

\infax[\ruledef{wfs-top}]
  {\Gamma\vdash\top\ \wf}
\end{minipage}\hfill
\begin{minipage}[t]{0.32\linewidth}
\infrule[\ruledef{wfs-single}]
  {\Gamma\vdash p\Rightarrow T}
  {\Gamma\vdash\sing p\ \wf}

\infrule[\ruledef{wfs-select}]
  {\Gamma\vdash p.a\Rightarrow S..R
   \andalso\Gamma\vdash S<:R}
  {\Gamma\vdash p.a\ \wf}
\end{minipage}\hfill
\begin{minipage}[t]{0.34\linewidth}
\infrule[\ruledef{wfs-fun}]
  {\Gamma\vdash T\ \wf
   \andalso\Gamma,x:T\vdash U\ \wf}
  {\Gamma\vdash(x:T)\to U\ \wf}

\infrule[\ruledef{wfs-pair}]
  {\Gamma\vdash T\ \wf
   \andalso\Gamma,x:T\vdash\mu\ \wf}
  {\Gamma\vdash\pairshape{x}{T}{a}{\mu}\ \wf}
\end{minipage}

\caption{Source syntax and static semantics: source syntax, path synthesis,
and well-formedness judgments.}
\Description{The complete source grammar, path-synthesis judgment, and all
well-formedness judgments.}
\label{fig:syntax-wf}
\label{fig:path-synthesis}
\label{fig:well-formedness}
\end{figure}

\clearpage
\begin{figure}[H]
\ContinuedFloat
\plateheader{Subcapturing}{\Gamma\vdash C<:D}
\noindent
\begin{minipage}[t]{0.31\linewidth}
\infax[\ruledef{sc-refl}]
  {\Gamma\vdash C<:C}

\infrule[\ruledef{sc-trans}]
  {\Gamma\vdash C<:D
   \andalso
   \Gamma\vdash D<:C'}
  {\Gamma\vdash C<:C'}

\infax[\ruledef{sc-empty}]
  {\Gamma\vdash\set{}<:C}

\infax[\ruledef{sc-union-left}]
  {\Gamma\vdash C<:C\cup D}
\end{minipage}\hfill
\begin{minipage}[t]{0.32\linewidth}
\infax[\ruledef{sc-union-right}]
  {\Gamma\vdash D<:C\cup D}

\infrule[\ruledef{sc-union-elim}]
  {\Gamma\vdash C<:C'
   \andalso
   \Gamma\vdash D<:C'}
  {\Gamma\vdash C\cup D<:C'}

\infrule[\ruledef{sc-path}]
  {\Gamma\vdash p\Rightarrow S\capt C}
  {\Gamma\vdash\set p<:C}

\infrule[\ruledef{sc-alias}]
  {\Gamma\vdash p\Rightarrow\sing q\capt C}
  {\Gamma\vdash\set q<:\set p}
\end{minipage}\hfill
\begin{minipage}[t]{0.34\linewidth}
\infrule[\ruledef{sc-fst-root}]
  {\Gamma\vdash p.1\Rightarrow T}
  {\Gamma\vdash\set{p.1}<:\set p}

\infrule[\ruledef{sc-sel-root}]
  {\Gamma\vdash p.a\Rightarrow T}
  {\Gamma\vdash\set{p.a}<:\set p}

\infrule[\ruledef{sc-select-lower}]
  {\Gamma\vdash p.a\Rightarrow L..U
   \andalso
   \Gamma\vdash L<:U}
  {\Gamma\vdash L<:p.a}

\infrule[\ruledef{sc-select-upper}]
  {\Gamma\vdash p.a\Rightarrow L..U
   \andalso
   \Gamma\vdash L<:U}
  {\Gamma\vdash p.a<:U}
\end{minipage}

\plateheader{Type subtyping}{\Gamma\vdash T<:U}
\noindent
\begin{minipage}[t]{0.30\linewidth}
\infax[\ruledef{st-refl}]
  {\Gamma\vdash T<:T}
\end{minipage}\hfill
\begin{minipage}[t]{0.31\linewidth}
\infrule[\ruledef{st-trans}]
  {\Gamma\vdash T<:U
   \andalso
   \Gamma\vdash U<:T'}
  {\Gamma\vdash T<:T'}
\end{minipage}\hfill
\begin{minipage}[t]{0.35\linewidth}
\infrule[\ruledef{st-capt}]
  {\Gamma\vdash C<:D
   \andalso
   \Gamma\vdash S<:R}
  {\Gamma\vdash S\capt C<:R\capt D}
\end{minipage}

\plateheader{Shape-type subtyping}{\Gamma\vdash S<:R}
\noindent
\begin{minipage}[t]{0.31\linewidth}
\infax[\ruledef{ss-refl}]
  {\Gamma\vdash S<:S}

\infrule[\ruledef{ss-trans}]
  {\Gamma\vdash S<:R
   \andalso
   \Gamma\vdash R<:S'}
  {\Gamma\vdash S<:S'}

\infax[\ruledef{ss-bot}]
  {\Gamma\vdash\bot<:S}

\infax[\ruledef{ss-top}]
  {\Gamma\vdash S<:\top}
\end{minipage}\hfill
\begin{minipage}[t]{0.32\linewidth}
\infrule[\ruledef{ss-widen}]
  {\Gamma\vdash p\Rightarrow S\capt C}
  {\Gamma\vdash\sing p<:S}
\infrule[\ruledef{ss-alias}]
  {\Gamma\vdash p\Rightarrow\sing q\capt C}
  {\Gamma\vdash\sing q<:\sing p}

\infrule[\ruledef{ss-select-lower}]
  {\Gamma\vdash p.a\Rightarrow S..R
   \andalso
   \Gamma\vdash S<:R}
  {\Gamma\vdash S<:p.a}
\end{minipage}\hfill
\begin{minipage}[t]{0.34\linewidth}
\infrule[\ruledef{ss-select-upper}]
  {\Gamma\vdash p.a\Rightarrow S..R
   \andalso
   \Gamma\vdash S<:R}
  {\Gamma\vdash p.a<:R}

\infrule[\ruledef{ss-fun}]
  {\Gamma\vdash T'<:T
   \andalso
   \Gamma,x:T'\vdash U<:U'}
  {\Gamma\vdash(x:T)\to U<:(x:T')\to U'}

\infrule[\ruledef{ss-pair}]
  {\Gamma\vdash T<:T'
   \andalso
   \Gamma,x:T\vdash\mu<:\mu'}
  {\Gamma\vdash
    \pairshape{x}{T}{a}{\mu}<:
    \pairshape{x}{T'}{a}{\mu'}}
\end{minipage}

\plateheader{Member-signature subtyping}{\Gamma\vdash\mu<:\nu}
\noindent
\begin{minipage}[t]{0.47\linewidth}
\infax[\ruledef{sm-refl}]
  {\Gamma\vdash\mu<:\mu}

\infrule[\ruledef{sm-trans}]
  {\Gamma\vdash\mu<:\nu
   \andalso
   \Gamma\vdash\nu<:\mu'}
  {\Gamma\vdash\mu<:\mu'}

\infrule[\ruledef{sm-term}]
  {\Gamma\vdash T<:U}
  {\Gamma\vdash T<:U}
\end{minipage}\hfill
\begin{minipage}[t]{0.51\linewidth}
\infrule[\ruledef{sm-type}]
  {\Gamma\vdash S'<:S
   \andalso
   \Gamma\vdash R<:R'
   \andalso
   \Gamma\vdash S<:R}
  {\Gamma\vdash S..R<:S'..R'}

\infrule[\ruledef{sm-capture}]
  {\Gamma\vdash C'<:C
   \andalso
   \Gamma\vdash D<:D'
   \andalso
   \Gamma\vdash C<:D}
  {\Gamma\vdash C..D<:C'..D'}
\end{minipage}

\caption{Static semantics (continued): subcapturing and all subtyping
judgments.}
\Description{The complete subcapturing, type, shape-type, and
member-signature subtyping judgments.}
\label{fig:subcapturing}
\label{fig:subtyping}
\label{fig:member-subtyping}
\end{figure}

\clearpage
\begin{figure}[H]
\ContinuedFloat

\plateheader{Term typing}{\typ{C}{\Gamma}{t}{T}}
\noindent
\begin{minipage}[t]{0.47\linewidth}
\infrule[\ruledef{t-path}]
  {\Gamma\vdash p\Rightarrow T}
  {\typ{\set p}{\Gamma}{p}{\sing p\capt\set p}}

\infrule[\ruledef{t-abs}]
  {\typ{C\cup\set x}{\Gamma,x:T}{t}{U}
   \\
   \Gamma\vdash T\ \wf
   \andalso
   \Gamma\vdash C\ \wf}
  {\typ{\set{}}{\Gamma}{\lambda(x:T).t}{((x:T)\to U)\capt C}}

\infrule[\ruledef{t-app}]
  {\typ{C_p}{\Gamma}{p}{((x:T)\to U)\capt C}
   \\
   \typ{C_q}{\Gamma}{q}{T}}
  {\typ{C_p\cup C_q}{\Gamma}{p\,q}{\subst{U}{q}{x}}}

\infax[\ruledef{t-pair}]
  {\begin{gathered}
   \set{};\,\Gamma\vdash\pairval{y}{a}{z}:\\[-2pt]
   \pairshape{x}{\sing y\capt\set y}{a}
     {\sing z\capt\set z}
     \capt(\set y\cup\set z)
   \end{gathered}}
\end{minipage}\hfill
\begin{minipage}[t]{0.51\linewidth}
\infrule[\ruledef{t-type-pair}]
  {\Gamma\vdash S\ \wf}
  {\typ{\set{}}{\Gamma}{\pairval{y}{a}{S}}{
   \pairshape{x}{\sing y\capt\set y}{a}
     {S..S}\capt\set y}}

\infrule[\ruledef{t-capture-pair}]
  {\Gamma\vdash C\ \wf}
  {\typ{\set{}}{\Gamma}{\pairval{y}{a}{C}}{
   \pairshape{x}{\sing y\capt\set y}{a}
     {C..C}\capt\set y}}

\infrule[\ruledef{t-let}]
  {\typ{C}{\Gamma}{t}{T}
   \andalso
   \typ{C}{\Gamma,x:T}{u}{U}
   \\
   \Gamma\vdash U\ \wf
   \andalso
   \Gamma\vdash C\ \wf}
  {\typ{C}{\Gamma}{\LET x=t\ \IN u}{U}}

\infrule[\ruledef{t-sub}]
  {\typ{C}{\Gamma}{t}{T}
   \andalso
   \Gamma\vdash T<:U
   \andalso
   \Gamma\vdash C<:D
   \\
   \Gamma\vdash U\ \wf
   \andalso
   \Gamma\vdash D\ \wf}
  {\typ{D}{\Gamma}{t}{U}}
\end{minipage}

\noindent
\begin{minipage}[t]{0.59\linewidth}
\plateheader{Runtime syntax}{}
\[
\begin{array}{rcl@{\qquad}l}
r &::=& \loc x\mid\typeref S\mid\capref C & \text{referent}\cr
\sigma &::=& \emptyset\mid\sigma\mathbin{\cdot}(x\mapsto v)
  & \text{store}\cr
k &::=& [\,]\mid
  (\LET x=[\,]\ \IN u)::k
  & \text{continuation}\cr
s &::=& \langle\sigma,k,t\rangle & \text{state}
\end{array}
\]
\[
\refdef z=\loc z\qquad
\refdef S=\typeref S\qquad
\refdef C=\capref C.
\]
\end{minipage}\hfill
\begin{minipage}[t]{0.38\linewidth}
\plateheader{Final state}{\mathsf{final}(s)}
\infax[\ruledef{f-location}]
  {\mathsf{final}(\langle\sigma,[\,],x\rangle)}

\infrule[\ruledef{f-value}]
  {v\ \mathsf{value}}
  {\mathsf{final}(\langle\sigma,[\,],v\rangle)}
\end{minipage}

\plateheader{Generalized path resolution}{\sigma\vdash p\resolve r}
\noindent
\begin{minipage}[t]{0.47\linewidth}
\infax[\ruledef{r-var}]
  {\sigma\vdash x\resolve\loc x}

\infrule[\ruledef{r-fst}]
  {\sigma\vdash p\resolve\loc z
   \andalso
   \sigma(z)=\pairval{y}{a}{\delta}}
  {\sigma\vdash p.1\resolve\loc y}
\end{minipage}\hfill
\begin{minipage}[t]{0.51\linewidth}
\infrule[\ruledef{r-sel}]
  {\sigma\vdash p\resolve\loc z
   \andalso
   \sigma(z)=\pairval{y}{a}{\delta}}
  {\sigma\vdash p.a\resolve\refdef\delta}

\infrule[\ruledef{r-miss}]
  {\sigma\vdash p\resolve\loc z
   \andalso
   \sigma(z)=\pairval{y}{b}{\delta}
   \\
   a\ne b
   \andalso
   \sigma\vdash y.a\resolve r}
  {\sigma\vdash p.a\resolve r}
\end{minipage}

\plateheader{Machine transition}
 {\langle\sigma,k,t\rangle\sstep\langle\sigma',k',t'\rangle}
\noindent
\begin{minipage}[t]{0.47\linewidth}
\infrule[\ruledef{e-app}]
  {\sigma\vdash p\resolve\loc z
   \andalso
   \sigma\vdash q\resolve\loc y
   \\
   \sigma(z)=\lambda(x:T).t}
  {\langle\sigma,k,p\,q\rangle\sstep
   \langle\sigma,k,\subst{t}{y}{x}\rangle}

\infrule[\ruledef{e-path}]
  {\sigma\vdash p\resolve\loc x
   \andalso p\ \mathsf{not\ a\ variable}}
  {\langle\sigma,k,p\rangle\sstep
   \langle\sigma,k,x\rangle}

\infax[\ruledef{e-let}]
  {\langle\sigma,k,\LET x=t\ \IN u\rangle
   \sstep
   \langle\sigma,(\LET x=[\,]\ \IN u)::k,t\rangle}
\end{minipage}\hfill
\begin{minipage}[t]{0.51\linewidth}
\infax[\ruledef{e-return}]
  {\langle\sigma,
      (\LET x=[\,]\ \IN u)::k,y\rangle
   \sstep
   \langle\sigma,k,\subst{u}{y}{x}\rangle}

\infrule[\ruledef{e-allocate}]
  {v\ \mathsf{value}
   \andalso \sigma(z)\ \mathsf{undefined}}
  {\langle\sigma,
      (\LET x=[\,]\ \IN u)::k,v\rangle
   \sstep\\
   \langle\sigma\mathbin{\cdot}(z\mapsto v),k,
     \subst{u}{z}{x}\rangle}
\end{minipage}

\caption{Static and dynamic semantics (continued): term typing, runtime
syntax, path resolution, and machine transitions.}
\Description{The complete term-typing judgment, runtime grammar,
final-state judgment, generalized path-resolution judgment, and CK-machine
transition relation.}
\label{fig:term-typing}
\label{fig:typing-runtime}
\label{fig:runtime}
\end{figure}

\endgroup
\clearpage

Path synthesis \(\Gamma\vdash p\Rightarrow\mu\) is precise and has no
subsumption. The capture set of a pair type is irrelevant to lookup. Matching
selection opens the dependent member with the receiver's first projection.
When labels differ, selection continues through the first component; this
represents a right-nested record without a separate record calculus.

There are four subtyping judgments. Subcapturing
\(\Gamma\vdash C<:D\) compares capture sets. Type and shape-type
subtyping are written \(\Gamma\vdash T<:U\) and
\(\Gamma\vdash S<:R\); the syntactic categories disambiguate them.
Member-signature subtyping uses the same notation. Interval
well-formedness and the type- and capture-set-selection subtyping rules carry
consistency premises. A capturing type may also be regarded as a
term signature. Since the grammar presents this inclusion without a
constructor, the premise and conclusion of the corresponding
well-formedness and subtyping rules have the same typography but belong to
different judgment families.

The premise use set \(C\cup\set x\) of the abstraction rule contains no free
occurrence of the new binder after abstraction. Subsumption and the term-path
root rules allow any path
rooted at \(x\), such as \(x.1.a\), to be widened to \(\set x\). Let
typing similarly requires the body's result type and use set to be
formed in the outer context; subsumption must eliminate the local root before
the rule applies.

We write
\(\mathsf{initial}(t)=\langle\emptyset,[\,],t\rangle\). A store is an
immutable sequence of allocated values. In particular, every runtime variable names
a store cell.  Allocation chooses a fresh location and transports the old
store, continuation, and term into the extended location scope.  The named
presentation of \ruleref{e-allocate} leaves this transport implicit.

Call a member signature bound-consistent in \(\Gamma\) when it is a
term signature, or when it is an interval whose lower bound is below its
upper bound in the corresponding subtype judgment.

\begin{lemma}[Bound consistency]
\label{lem:bound-consistency}
Every well-formed member signature is bound-consistent. If
\(\Gamma\vdash\mu<:\nu\) and \(\mu\) is bound-consistent, then \(\nu\) is
bound-consistent.
\end{lemma}

The well-formedness result follows immediately from the interval rules.  For
signature subtyping, the type-interval case chains
\(S'<:S<:R<:R'\), and the capture-set-interval case chains
\(C'<:C<:D<:D'\); reflexivity and transitivity preserve the property.
The mechanization does not yet include a full context-formation and typing
regularity theorem.  That stronger result requires the opening and path
substitution metatheory needed to recover well-formed synthesized signatures
from arbitrary typing derivations.

\begin{lemma}[Resolution determinism]
\label{lem:resolution-determinism}
If \(\sigma\vdash p\resolve r\) and
\(\sigma\vdash p\resolve r'\), then \(r=r'\).
\end{lemma}

The proof follows the first resolution derivation.  In the two selection
cases, immutability gives a unique store binding, and the induction
hypothesis determines the tail referent.  Generalized resolution is used in
the static and semantic arguments; the machine steps only when its result is
a location.

\subsection{Capture-set members and root contraction}

A capture-set definition packages a capture set behind a stable path. For
example, after allocating
\[
  z=\pairval{y}{a}{D},
\]
path synthesis exposes the exact interval
\(z.a\Rightarrow D..D\). Subcapturing may subsequently forget precision by
placing \(z.a\) between wider declared bounds. Capture-set selection therefore
provides capture-set abstraction without adding a separate binder form or
explicit capture polymorphism to terms.

The root rules apply only to singleton capture sets of term paths.
They express that using \(p.1\) or \(p.a\) requires access to the root
capability \(p\). An abstract capture-set selection has no such ownership
interpretation. This restriction is visible in the type of a concrete
capture-set pair: its capture set is \(\set y\), independent of the definition
\(D\), which is static metadata. If one added \(p.a<:\set p\), the exact lower
bound \(D<:p.a\), together with \(\set p<:\set y\), would derive
\(D<:\set y\) for an arbitrary definition. This would impose an ownership
constraint on static metadata that the value-introduction rule does not
establish.

The two syntactic occurrences of paths in capture sets also remain distinct.
The singleton \(\set p\) denotes the capability named by a term path; the
selection \(p.a\) denotes the capture set stored in a capture-set member. This
distinction is used by both path synthesis and store-indexed subcapturing.

\subsection{The general pair rule}

Rule \ruleref{ss-pair} compares dependent members under the source type of
the first component.  It supports ordinary covariance of the first component
without requiring the member signature to be independent of the binder.  Two
representative instances are
\[
\begin{array}{c}
 x:\top\vdash \bot<:\sing x
 \qquad
 x:\top\vdash \sing x<:\top
 \\[-1pt]
 \hline
 \pairshape{x}{\top}{a}{\sing x..\sing x}
 <: \pairshape{x}{\top}{a}{\bot..\top}
\end{array}
\qquad
\begin{array}{c}
 x:\top\vdash \set{}<:\set x
 \qquad
 x:\top\vdash \set x<:\set x
 \\[-1pt]
 \hline
 \pairshape{x}{\top}{a}{\set x..\set x}
 <: \pairshape{x}{\top}{a}{\set{}..\set x}.
\end{array}
\]
The mechanization exercises both derivations in closed executions that
allocate the coerced pair, thereby forcing the semantic pair action rather
than merely constructing a source derivation. A separate capture-set-selection
execution forces both the lower- and upper-bound selection coercions.  The
examples are useful because the member premises genuinely mention \(x\), even
though the first-component types on both sides are \(\top\).

\section{Store-indexed interpretation of capturing types}
\label{sec:semantics}

The preservation argument translates source subtyping derivations into
store-indexed coercions. Coercion application maps runtime typing at a source
type to runtime typing at its target type, so preservation does not invert a
derivation ending in transitivity. The interpretation also retains the
capture set assigned to every stored value and the use set of every running
term.

\subsection{Runtime path equality and conversion}

The store equates paths that resolve to the same referent.  We write
\(p\equiv_\sigma q\) for the least equivalence and path congruence generated
by co-resolution.

\begin{figure}[H]
\begingroup
\rulesetup
\plateheader{Runtime path equality}{p\equiv_\sigma q}
\noindent
\begin{minipage}[t]{0.47\linewidth}
\infax[\ruledef{peq-refl}]
  {p\equiv_\sigma p}

\infrule[\ruledef{peq-symm}]
  {p\equiv_\sigma q}
  {q\equiv_\sigma p}

\infrule[\ruledef{peq-trans}]
  {p\equiv_\sigma q\andalso q\equiv_\sigma p'}
  {p\equiv_\sigma p'}
\end{minipage}\hfill
\begin{minipage}[t]{0.51\linewidth}
\infrule[\ruledef{peq-coresolve}]
  {\sigma\vdash p\resolve r
   \andalso
   \sigma\vdash q\resolve r}
  {p\equiv_\sigma q}

\infrule[\ruledef{peq-fst}]
  {p\equiv_\sigma q}
  {p.1\equiv_\sigma q.1}

\infrule[\ruledef{peq-sel}]
  {p\equiv_\sigma q}
  {p.a\equiv_\sigma q.a}
\end{minipage}
\endgroup
\caption{Store-induced equality of paths.}
\Description{All rules defining runtime path equality.}
\label{fig:runtime-equality}
\end{figure}

Runtime conversion \(X\simeq_\sigma Y\) is defined for capture sets,
types, shape types, and member signatures. For each category it is generated
inductively by reflexivity, replacement, and congruence for every constructor.
Replacement has the schema
\[
  \frac{p\equiv_\sigma q}
       {\subst{X}{p}{x}\simeq_\sigma\subst{X}{q}{x}}.
\]
Reverse replacements use symmetry of path equality; sequences of conversions
are composed by the semantic coercion judgments below. Conversion preserves
the outer constructor and the form of a member signature. Beneath a binder,
ambient equalities are weakened and the bound variable is related only to
itself; we write the resulting conversion as
\(X\simeq^x_\sigma Y\).  Opening it with runtime-equal paths yields ordinary
conversion.

\subsection{Annotated stores and semantic judgments}

An annotated store \(\Omega\) has the same spine as \(\sigma\) and records
beside each stored value the capture set assigned by its introduction rule:
\[
  \Omega ::= [\,] \mid \Omega\mathbin{\cdot}(v@Q).
\]
We call \(Q\) the \emph{assigned capture set} of \(v\), and write
\(\Omega(x)=v@Q\) for lookup. All locations in \(v\) and \(Q\)
refer to earlier cells, as they do in the underlying store. The judgment
\(\Omega\valid\sigma\), defined after value typing, states
that each annotation is justified by the introduction rule for the
corresponding stored value.

A valuation for a context \(\Gamma\) in a store \(\sigma\) is a total function
\[
  \rho : \mathsf{Fin}(|\Gamma|)\longrightarrow
         \mathsf{Fin}(|\sigma|),
\]
where \(|\Gamma|\) and \(|\sigma|\) are the numbers of context bindings and
store cells. The actions \(\rho(p)\), \(\rho(t)\), and \(\rho(X)\) replace
every free variable in a path, term, type, shape type, capture set, or member
signature with the location selected by \(\rho\), while leaving a displayed
binder in place. Valuations need not be injective.

The semantic judgments are
\[
\begin{array}{rcll}
 \rho\models^\Omega_\sigma\Gamma
   &&& \text{environment satisfaction},\\
 \Omega;\sigma\vdash C\semleq D
   &&& \text{store-indexed subcapturing},\\
 \Omega;\sigma\vdash T\tcoerce U
   &&& \text{type coercion},\\
 \Omega;\sigma\vdash S\scoerce R
   &&& \text{shape-type coercion},\\
 \Omega;\sigma\vdash\mu\coerce\nu
   &&& \text{member coercion},\\
 \Omega;\sigma;x:T\vdash U\rightsquigarrow_{\mathsf d}U'
   &&& \text{deferred function-result coercion},\\
 \mathsf{Body}_{\Omega;\sigma}(\typ{C}{x:T}{t}{U})
   &&& \text{suspended body typing},\\
 \mathsf{Member}_{\Omega;\sigma}(x:T\vdash\mu<:\nu)
   &&& \text{suspended member premise},\\
 \Omega;\sigma\models x:T@Q
   &&& \text{location typing with assigned capture set }Q,\\
 \Omega;\sigma\models r:\mu
   &&& \text{referent typing},\\
 \Omega;\sigma\vdash_{\mathsf v}v:T@Q
   &&& \text{value typing}.
\end{array}
\]
These judgments are mutually defined. Annotated store typing
\(\Omega\valid\sigma\) is subsequently defined by empty-store and allocation
rules; allocation requires value typing for the new entry.

\paragraph{Proof structure.}
Once \(\rho\models^\Omega_\sigma\Gamma\) holds, source subcapturing and
subtyping derivations compile to the corresponding store-indexed relations.
Coercion action then transports referent typing:
\[
\begin{array}{c@{\qquad}c@{\qquad}c}
 \Gamma\vdash\mu<:\nu
 &\Longrightarrow&
 \Omega;\sigma\vdash\rho(\mu)\coerce\rho(\nu),
 \\[3pt]
 \Omega;\sigma\models r:\rho(\mu)
 \quad
 \Omega;\sigma\vdash\rho(\mu)\coerce\rho(\nu)
 &\Longrightarrow&
 \Omega;\sigma\models r:\rho(\nu).
\end{array}
\]
The annotation \(Q\) is unchanged by this action, while the derivation from
\(Q\) to the capture set of the assigned type is composed with subcapturing.
For dependent pairs, compilation retains the member premise under its binder;
action instantiates that premise only after lookup supplies the stored first
component. The following figures define the relations used by these two
steps.

\clearpage
\begingroup
\setlength{\abovedisplayskip}{5pt plus 1pt minus 1pt}
\setlength{\belowdisplayskip}{5pt plus 1pt minus 1pt}
\setlength{\intextsep}{5pt}
\captionsetup{skip=6pt}
\rulesetup

\begin{figure}[H]
\plateheader{Store-indexed subcapturing}{\Omega;\sigma\vdash C\semleq D}
\noindent
\begin{minipage}[t]{0.31\linewidth}
\infrule[\ruledef{cap-source}]
  {\rho\models^\Omega_\sigma\Gamma
   \\
   \Gamma\vdash C<:D}
  {\Omega;\sigma\vdash\rho(C)\semleq\rho(D)}

\infax[\ruledef{cap-refl}]
  {\Omega;\sigma\vdash C\semleq C}

\infrule[\ruledef{cap-trans}]
  {\Omega;\sigma\vdash C\semleq D
   \\
   \Omega;\sigma\vdash D\semleq C'}
  {\Omega;\sigma\vdash C\semleq C'}

\infrule[\ruledef{cap-conv}]
  {C\simeq_\sigma D}
  {\Omega;\sigma\vdash C\semleq D}

\infax[\ruledef{cap-empty}]
  {\Omega;\sigma\vdash\set{}\semleq C}
\end{minipage}\hfill
\begin{minipage}[t]{0.31\linewidth}
\infax[\ruledef{cap-union-left}]
  {\Omega;\sigma\vdash C\semleq C\cup D}

\infax[\ruledef{cap-union-right}]
  {\Omega;\sigma\vdash D\semleq C\cup D}

\infrule[\ruledef{cap-union-elim}]
  {\Omega;\sigma\vdash C\semleq C'
   \\
   \Omega;\sigma\vdash D\semleq C'}
  {\Omega;\sigma\vdash C\cup D\semleq C'}

\infrule[\ruledef{cap-alias}]
  {\sigma\vdash p\resolve\loc x
   \\
   \sigma\vdash q\resolve\loc x}
  {\Omega;\sigma\vdash\set q\semleq\set p}

\infrule[\ruledef{cap-fold}]
  {\sigma\vdash p\resolve\loc x
   \\
   \Omega(x)=v@Q}
  {\Omega;\sigma\vdash Q\semleq\set p}
\end{minipage}\hfill
\begin{minipage}[t]{0.34\linewidth}
\infrule[\ruledef{cap-fst-root}]
  {\sigma\vdash p.1\resolve\loc x}
  {\Omega;\sigma\vdash\set{p.1}\semleq\set p}

\infrule[\ruledef{cap-sel-root}]
  {\sigma\vdash p.a\resolve\loc x}
  {\Omega;\sigma\vdash\set{p.a}\semleq\set p}

\infrule[\ruledef{cap-select-lower}]
  {\sigma\vdash p.a\resolve\capref C
   \\
   \Omega;\sigma\vdash L\semleq C}
  {\Omega;\sigma\vdash L\semleq p.a}

\infrule[\ruledef{cap-select-upper}]
  {\sigma\vdash p.a\resolve\capref C
   \\
   \Omega;\sigma\vdash C\semleq U}
  {\Omega;\sigma\vdash p.a\semleq U}
\end{minipage}

\caption{The complete store-indexed subcapturing relation.}
\Description{All rules of store-indexed subcapturing.}
\label{fig:semantic-subcapturing}
\end{figure}

\begin{figure}[H]
\ContinuedFloat
\plateheader{Type coercion}{\Omega;\sigma\vdash T\tcoerce U}
\noindent
\begin{minipage}[t]{0.22\linewidth}
\infax[\ruledef{tc-refl}]
  {\Omega;\sigma\vdash T\tcoerce T}
\end{minipage}\hfill
\begin{minipage}[t]{0.27\linewidth}
\infrule[\ruledef{tc-trans}]
  {\Omega;\sigma\vdash T\tcoerce U
   \\
   \Omega;\sigma\vdash U\tcoerce T'}
  {\Omega;\sigma\vdash T\tcoerce T'}
\end{minipage}\hfill
\begin{minipage}[t]{0.20\linewidth}
\infrule[\ruledef{tc-conv}]
  {T\simeq_\sigma U}
  {\Omega;\sigma\vdash T\tcoerce U}
\end{minipage}\hfill
\begin{minipage}[t]{0.27\linewidth}
\infrule[\ruledef{tc-capt}]
  {\Omega;\sigma\vdash C\semleq D
   \\
   \Omega;\sigma\vdash S\scoerce R}
  {\Omega;\sigma\vdash S\capt C\tcoerce R\capt D}
\end{minipage}

\caption{The complete coercion relation for capturing types.}
\Description{All rules for coercions between capturing types.}
\label{fig:type-coercion}
\end{figure}

\begin{figure}[H]
\ContinuedFloat
\plateheader{Shape-type coercion}{\Omega;\sigma\vdash S\scoerce R}
\noindent
\begin{minipage}[t]{0.30\linewidth}
\infax[\ruledef{shc-refl}]
  {\Omega;\sigma\vdash S\scoerce S}

\infrule[\ruledef{shc-trans}]
  {\Omega;\sigma\vdash S\scoerce R
   \\
   \Omega;\sigma\vdash R\scoerce S'}
  {\Omega;\sigma\vdash S\scoerce S'}

\infrule[\ruledef{shc-conv}]
  {S\simeq_\sigma R}
  {\Omega;\sigma\vdash S\scoerce R}

\infax[\ruledef{shc-bot}]
  {\Omega;\sigma\vdash\bot\scoerce S}

\infax[\ruledef{shc-top}]
  {\Omega;\sigma\vdash S\scoerce\top}
\end{minipage}\hfill
\begin{minipage}[t]{0.30\linewidth}
\infrule[\ruledef{shc-widen}]
  {\sigma\vdash p\resolve\loc x
   \\
   \Omega;\sigma\models x:S\capt C@Q}
  {\Omega;\sigma\vdash\sing p\scoerce S}

\infrule[\ruledef{shc-alias}]
  {\sigma\vdash p\resolve\loc x
   \\
   \sigma\vdash q\resolve\loc x}
  {\Omega;\sigma\vdash\sing q\scoerce\sing p}

\infrule[\ruledef{shc-select-lower}]
  {\sigma\vdash p.a\resolve\typeref R
   \\
   \Omega;\sigma\vdash L\scoerce R}
  {\Omega;\sigma\vdash L\scoerce p.a}

\infrule[\ruledef{shc-select-upper}]
  {\sigma\vdash p.a\resolve\typeref R
   \\
   \Omega;\sigma\vdash R\scoerce U}
  {\Omega;\sigma\vdash p.a\scoerce U}
\end{minipage}\hfill
\begin{minipage}[t]{0.36\linewidth}
\infrule[\ruledef{shc-fun}]
  {\Omega;\sigma\vdash T'\tcoerce T
   \\
   \Omega;\sigma;x:T'\vdash
      U\rightsquigarrow_{\mathsf d}U'}
  {\Omega;\sigma\vdash
    (x:T)\to U\scoerce(x:T')\to U'}

\infrule[\ruledef{shc-pair}]
  {\Omega;\sigma\vdash T\tcoerce T'
   \\
   \mathsf{Member}_{\Omega;\sigma}
      (x:T\vdash\mu<:\mu')}
  {\Omega;\sigma\vdash
    \pairshape{x}{T}{a}{\mu}\scoerce
    \pairshape{x}{T'}{a}{\mu'}}
\end{minipage}

\caption{The complete coercion relation for shape types.}
\Description{All rules for shape-type coercion.}
\label{fig:shape-coercion}
\end{figure}

\begin{figure}[H]
\ContinuedFloat
\plateheader{Member coercion}
  {\Omega;\sigma\vdash\mu\coerce\nu}
\noindent
\begin{minipage}[t]{0.31\linewidth}
\infax[\ruledef{mc-refl}]
  {\Omega;\sigma\vdash\mu\coerce\mu}

\infrule[\ruledef{mc-trans}]
  {\Omega;\sigma\vdash\mu\coerce\nu
   \\
   \Omega;\sigma\vdash\nu\coerce\mu'}
  {\Omega;\sigma\vdash\mu\coerce\mu'}
\end{minipage}\hfill
\begin{minipage}[t]{0.31\linewidth}
\infrule[\ruledef{mc-conv}]
  {\mu\simeq_\sigma\nu}
  {\Omega;\sigma\vdash\mu\coerce\nu}

\infrule[\ruledef{mc-term}]
  {\Omega;\sigma\vdash T\tcoerce U}
  {\Omega;\sigma\vdash T\coerce U}
\end{minipage}\hfill
\begin{minipage}[t]{0.34\linewidth}
\infrule[\ruledef{mc-type}]
  {\Omega;\sigma\vdash L'\scoerce L
   \andalso
   \Omega;\sigma\vdash U\scoerce U'}
  {\Omega;\sigma\vdash L..U\coerce L'..U'}

\infrule[\ruledef{mc-capture}]
  {\Omega;\sigma\vdash C'\semleq C
   \andalso
   \Omega;\sigma\vdash D\semleq D'}
  {\Omega;\sigma\vdash C..D\coerce C'..D'}
\end{minipage}

\caption{The complete member coercion relation.}
\Description{All rules for member coercion.}
\label{fig:member-coercions}
\end{figure}

\begin{figure}[H]
\ContinuedFloat
\plateheader{Suspended body}
 {\mathsf{Body}_{\Omega;\sigma}(\typ{C}{x:T}{t}{U})}
\infrule[\ruledef{body-source}]
  {\rho\models^\Omega_\sigma\Gamma
   \andalso
   \typ{C}{\Gamma,x:T_0}{t}{U}}
  {\mathsf{Body}_{\Omega;\sigma}
    (\typ{\rho(C)}{x:\rho(T_0)}{\rho(t)}{\rho(U)})}

\plateheader{Suspended member premise}
 {\mathsf{Member}_{\Omega;\sigma}(x:T\vdash\mu<:\nu)}
\infrule[\ruledef{member-source}]
  {\rho\models^\Omega_\sigma\Gamma
   \andalso
   \Gamma,x:T_0\vdash\mu<:\nu}
  {\mathsf{Member}_{\Omega;\sigma}
    (x:\rho(T_0)\vdash\rho(\mu)<:\rho(\nu))}

\plateheader{Deferred function-result coercion}
 {\Omega;\sigma;x:T\vdash U\rightsquigarrow_{\mathsf d}U'}
\noindent
\begin{minipage}[t]{0.47\linewidth}
\infax[\ruledef{dc-refl}]
  {\Omega;\sigma;x:T\vdash U\rightsquigarrow_{\mathsf d}U}

\infrule[\ruledef{dc-trans}]
  {\Omega;\sigma;x:T\vdash U\rightsquigarrow_{\mathsf d}U'
   \\
   \Omega;\sigma;x:T\vdash U'\rightsquigarrow_{\mathsf d}U''}
  {\Omega;\sigma;x:T\vdash U\rightsquigarrow_{\mathsf d}U''}

\infrule[\ruledef{dc-conv}]
  {U\simeq^x_\sigma U'}
  {\Omega;\sigma;x:T\vdash U\rightsquigarrow_{\mathsf d}U'}
\end{minipage}\hfill
\begin{minipage}[t]{0.51\linewidth}
\infrule[\ruledef{dc-narrow}]
  {\Omega;\sigma\vdash T'\tcoerce T
   \\
   \Omega;\sigma;x:T\vdash U\rightsquigarrow_{\mathsf d}U'}
  {\Omega;\sigma;x:T'\vdash U\rightsquigarrow_{\mathsf d}U'}

\infrule[\ruledef{dc-source}]
  {\rho\models^\Omega_\sigma\Gamma
   \andalso
   \Gamma,x:T\vdash U<:U'}
  {\Omega;\sigma;x:\rho(T)\vdash
    \rho(U)\rightsquigarrow_{\mathsf d}\rho(U')}
\end{minipage}

\caption{Suspended source premises and the complete deferred
function-result coercion.}
\Description{The source rules for suspended bodies and dependent members,
and all rules for coercions suspended beneath a function binder.}
\label{fig:deferred-coercions}
\end{figure}

\endgroup

Rule \ruleref{cap-source} embeds a source subcapturing derivation together
with the environment that justifies its path premises. The other constructors
provide the properties needed by coercion action and reduction. In
particular, \ruleref{cap-fold} states that a path naming a stored value
subsumes its assigned capture set. There is no inverse rule.
Abstract capture-set selections are related only through their stored interval
bounds.

Function coercions retain their codomain premise until an argument location
is available.  Pair coercions retain the member premise until the stored
first component is available. Type and capture-set intervals are contravariant
in the lower bound and covariant in the upper bound. In the source-premise
rules, \(\rho\) acts on variables from \(\Gamma\); the displayed
\(x\) remains bound.

\subsection{Value typing and annotated stores}

Environment satisfaction, location typing, referent typing, and value typing
are shown in the next two figures. Location typing carries an explicit assigned
capture set \(Q\). Every location rule retrieves the same \(Q\) from
\(\Omega\) and proves \(Q\semleq C\), where \(C\) is the capture set of the
assigned type.

\begingroup
\setlength{\abovedisplayskip}{5pt plus 1pt minus 1pt}
\setlength{\belowdisplayskip}{5pt plus 1pt minus 1pt}
\setlength{\intextsep}{5pt}
\captionsetup{skip=6pt}
\rulesetup

\begin{figure}[H]
\plateheader{Environment satisfaction}{\rho\models^\Omega_\sigma\Gamma}
\infrule[\ruledef{sem-env}]
  {\forall x:\mathsf{Fin}(|\Gamma|).\;
   \exists Q.\;
   \Omega;\sigma\models\rho(x):\rho(\Gamma(x))@Q}
  {\rho\models^\Omega_\sigma\Gamma}

\plateheader{Referent typing}{\Omega;\sigma\models r:\mu}
\noindent
\begin{minipage}[t]{0.31\linewidth}
\infrule[\ruledef{sem-ref-loc}]
  {\Omega;\sigma\models x:T@Q}
  {\Omega;\sigma\models\loc x:T}
\end{minipage}\hfill
\begin{minipage}[t]{0.32\linewidth}
\infrule[\ruledef{sem-ref-type}]
  {\Omega;\sigma\vdash L\scoerce R
   \\
   \Omega;\sigma\vdash R\scoerce U}
  {\Omega;\sigma\models\typeref R:L..U}
\end{minipage}\hfill
\begin{minipage}[t]{0.34\linewidth}
\infrule[\ruledef{sem-ref-capture}]
  {\Omega;\sigma\vdash L\semleq C
   \\
   \Omega;\sigma\vdash C\semleq U}
  {\Omega;\sigma\models\capref C:L..U}
\end{minipage}

\plateheader{Location typing}{\Omega;\sigma\models x:T@Q}
\noindent
\begin{minipage}[t]{0.47\linewidth}
\infrule[\ruledef{sem-loc-top}]
  {\Omega(x)=v@Q
   \andalso
   \Omega;\sigma\vdash Q\semleq C}
  {\Omega;\sigma\models x:\top\capt C@Q}

\infrule[\ruledef{sem-loc-fun}]
  {\Omega(x)=\lambda(z:T_0).t@Q
   \\
   \mathsf{Body}_{\Omega;\sigma}
     (\typ{Q\cup\set z}{z:T_0}{t}{U_0})
   \\
   \Omega;\sigma\vdash T\tcoerce T_0
   \andalso
   \Omega;\sigma;z:T\vdash U_0\rightsquigarrow_{\mathsf d}U
   \\
   \Omega;\sigma\vdash Q\semleq C}
  {\Omega;\sigma\models
    x:((z:T)\to U)\capt C@Q}
\end{minipage}\hfill
\begin{minipage}[t]{0.51\linewidth}
\infrule[\ruledef{sem-loc-pair}]
  {\Omega(x)=\pairval{y}{a}{\delta}@Q
   \\
   \Omega;\sigma\models y:T@Q_y
   \\
   \Omega;\sigma\models
      \refdef\delta:\subst{\mu}{y}{z}
   \\
   \Omega;\sigma\vdash Q\semleq C}
  {\Omega;\sigma\models
    x:\pairshape{z}{T}{a}{\mu}\capt C@Q}

\infrule[\ruledef{sem-loc-single}]
  {\Omega(x)=v@Q
   \andalso
   \sigma\vdash p\resolve\loc x
   \\
   \Omega;\sigma\vdash Q\semleq C}
  {\Omega;\sigma\models x:\sing p\capt C@Q}

\infrule[\ruledef{sem-loc-select}]
  {\Omega(x)=v@Q
   \andalso
   \sigma\vdash p.a\resolve\typeref R
   \\
   \Omega;\sigma\models x:R\capt D@Q
   \andalso
   \Omega;\sigma\vdash Q\semleq C}
  {\Omega;\sigma\models x:p.a\capt C@Q}
\end{minipage}

\caption{Environment satisfaction, referent typing, and location typing.}
\Description{All rules for semantic environments, generalized referents,
and locations.}
\label{fig:store-typing-relations}
\end{figure}

\endgroup

The introduction types of pair values are abbreviated as follows:
\[
\begin{array}{rcl}
I_{\mathsf{term}}(y,a,z)
 &=& \pairshape{x}{\sing y\capt\set y}{a}
    {\sing z\capt\set z}
   \capt(\set y\cup\set z),
\\[2pt]
I_{\mathsf{type}}(y,a,R)
 &=& \pairshape{x}{\sing y\capt\set y}{a}
    {R..R}
   \capt\set y,
\\[2pt]
I_{\mathsf{capture}}(y,a,C)
 &=& \pairshape{x}{\sing y\capt\set y}{a}
    {C..C}
   \capt\set y.
\end{array}
\]
As above, weakening beneath the pair binder is implicit.

\begingroup
\setlength{\abovedisplayskip}{5pt plus 1pt minus 1pt}
\setlength{\belowdisplayskip}{5pt plus 1pt minus 1pt}
\setlength{\intextsep}{5pt}
\captionsetup{skip=6pt}
\rulesetup

\begin{figure}[H]
\plateheader{Value typing}
 {\Omega;\sigma\vdash_{\mathsf v}v:T@Q}
\noindent
\begin{minipage}[t]{0.47\linewidth}
\infrule[\ruledef{ve-abs}]
  {\mathsf{Body}_{\Omega;\sigma}
     (\typ{Q\cup\set x}{x:T_0}{t}{U_0})
   \\
   \Omega;\sigma\vdash
      ((x:T_0)\to U_0)\capt Q\tcoerce T}
  {\Omega;\sigma\vdash_{\mathsf v}
     \lambda(x:T_0).t:T@Q}

\infrule[\ruledef{ve-term-pair}]
  {\Omega;\sigma\vdash I_{\mathsf{term}}(y,a,z)\tcoerce T}
  {\Omega;\sigma\vdash_{\mathsf v}
     \pairval{y}{a}{z}:T@(\set y\cup\set z)}
\end{minipage}\hfill
\begin{minipage}[t]{0.51\linewidth}
\infrule[\ruledef{ve-type-pair}]
  {\Omega;\sigma\vdash I_{\mathsf{type}}(y,a,R)\tcoerce T}
  {\Omega;\sigma\vdash_{\mathsf v}
     \pairval{y}{a}{R}:T@\set y}

\infrule[\ruledef{ve-capture-pair}]
  {\Omega;\sigma\vdash I_{\mathsf{capture}}(y,a,C)\tcoerce T}
  {\Omega;\sigma\vdash_{\mathsf v}
     \pairval{y}{a}{C}:T@\set y}
\end{minipage}

\plateheader{Annotated store typing}{\Omega\valid\sigma}
\noindent
\begin{minipage}[t]{0.26\linewidth}
\infax[\ruledef{store-empty}]
  {[\,]\valid\emptyset}
\end{minipage}\hfill
\begin{minipage}[t]{0.68\linewidth}
\infrule[\ruledef{store-allocate}]
  {\Omega\valid\sigma
   \andalso
   \Omega;\sigma\vdash_{\mathsf v}v:T@Q
   \andalso v\ \mathsf{value}}
  {\Omega\mathbin{\cdot}(v@Q)\valid
   \sigma\mathbin{\cdot}(z\mapsto v)}
\end{minipage}

\caption{Value typing and annotated store typing.}
\Description{All rules for value typing and annotated store typing.}
\label{fig:value-store-typing}
\end{figure}

\endgroup

The assigned capture set \(Q\) following \(@\) is fixed by the introduction
rule. It need not be minimal: the body derivation in \ruleref{ve-abs}, for
example, may already contain subsumption. The coercion derivation records any further subsumption above
the value constructor. For every shape type \(S\),
\[
  \Omega;\sigma\vdash_{\mathsf v}v:S\capt C@Q
  \quad\Longrightarrow\quad
  \Omega;\sigma\vdash Q\semleq C.
\]
Annotated store typing ties each annotation to value typing for the same store entry.
Because stores are append-only, \ruleref{store-allocate} depends only on the
previous annotated store. Allocation preserves all earlier typing derivations by
weakening locations uniformly.

\subsection{Compilation and coercion action}

Typed resolution and coercion compilation are stated under a satisfied
source environment:
\[
\begin{array}{rcl}
\Gamma\vdash p\Rightarrow\mu
 &\Longrightarrow&
 \exists r.\;\sigma\vdash\rho(p)\resolve r
       \ \land\ \Omega;\sigma\models r:\rho(\mu),
\\[3pt]
\Gamma\vdash C<:D
 &\Longrightarrow&
 \Omega;\sigma\vdash\rho(C)\semleq\rho(D),
\\
\Gamma\vdash T<:U
 &\Longrightarrow&
 \Omega;\sigma\vdash\rho(T)\tcoerce\rho(U),
\\
\Gamma\vdash S<:R
 &\Longrightarrow&
 \Omega;\sigma\vdash\rho(S)\scoerce\rho(R),
\\
\Gamma\vdash\mu<:\nu
 &\Longrightarrow&
 \Omega;\sigma\vdash\rho(\mu)\coerce\rho(\nu).
\end{array}
\]
Each construction follows its source derivation. Transitivity builds an
explicit composition.  Function subtyping produces \ruleref{shc-fun} with a
deferred result coercion. Pair subtyping produces \ruleref{shc-pair} with a
suspended member premise. Neither case requires transitivity inversion.

\begin{lemma}[Coercion action]
\label{lem:coercion-action}
If \(\Omega;\sigma\models r:\mu\) and
\(\Omega;\sigma\vdash\mu\coerce\nu\), then
\(\Omega;\sigma\models r:\nu\).  In particular,
\[
 \Omega;\sigma\models x:T@Q
 \quad\land\quad
 \Omega;\sigma\vdash T\tcoerce U
 \quad\Longrightarrow\quad
 \Omega;\sigma\models x:U@Q.
\]
\end{lemma}

Identity returns its argument, composition applies its components in order,
and runtime conversion transports typing derivations along store-induced path
equality. Interval action composes the coercions on their bounds. For
capturing types, \ruleref{tc-capt} composes the existing derivation
\(Q\semleq C\) with its premise \(C\semleq D\). The shape-type component
changes the assigned shape type while retaining the same \(Q\). This step preserves the
returned-value capture-set bound through subsumption.

\subsection{Action for a dependent pair}
\label{sec:pair-action}

Suppose \(z\) has pair typing with assigned capture set \(Q\):
\[
 \Omega;\sigma\models
 z:\pairshape{x}{T}{a}{\mu}\capt C@Q,
 \qquad
 \Omega;\sigma\vdash Q\semleq C.
\]
Consider a coercion to a pair type with first-component type \(T'\), member
signature \(\mu'\), and capture set \(D\):
\[
 \Omega;\sigma\vdash
 \pairshape{x}{T}{a}{\mu}\capt C
 \tcoerce
 \pairshape{x}{T'}{a}{\mu'}\capt D.
\]
Inverting \ruleref{sem-loc-pair} exposes a first location \(y\), a member
referent \(r\), and the typing facts
\[
 \Omega;\sigma\models y:T@Q',
 \qquad
 \Omega;\sigma\models r:\subst{\mu}{y}{x}.
\]
The shape-type component of a compiled pair coercion contains
\[
 \Omega;\sigma\vdash T\tcoerce T',
 \qquad
 \mathsf{Member}_{\Omega;\sigma}(x:T\vdash\mu<:\mu').
\]
Coercion action first derives \(\Omega;\sigma\models y:T'@Q'\). It then
instantiates the suspended member premise with the same location \(y\),
using its typing at the source first-component type:
\[
 \Omega;\sigma\vdash
 \subst{\mu}{y}{x}\coerce\subst{\mu'}{y}{x}.
\]
Applying this coercion to the member's typing derives the target signature.
The outer pair retains the assigned capture set \(Q\); the enclosing
coercion's subcapturing premise composes \(Q\semleq C\semleq D\). This proves the general pair
rule for term, type, and capture-set members without requiring the dependent
member to ignore \(x\).

Coercion action on a member creates a new coercion after action on the outer
pair has started, so structural recursion on the original coercion is
insufficient. Define the allocation stratum of a location as one plus the
number of cells allocated before it, and that of a shape-type or capture-set
referent as zero. Thus newer locations have greater
strata. Pair cells refer only to older locations, so action on either stored
component decreases the stratum. Every other recursive call traverses a
strict subderivation and decreases the coercion tree size. The lexicographic
pair of allocation stratum and coercion tree size is therefore well founded.

\section{Runtime typing invariant}
\label{sec:invariant}

The machine invariant retains result types and use sets together. Source
subsumption is retained as a type-coercion derivation and a store-indexed
subcapturing derivation at the outer runtime constructor. The judgment
\(\semtyp{C}{\Omega;\sigma}{t}{T}\) assigns the running term its
current type and use set. A continuation judgment
\[
 \Omega;\sigma\vdash k:(U,D)\Rightarrow(T,C)
\]
says that a continuation accepting a current result of type \(U\) and use
set \(D\) produces final type \(T\) with use set \(C\).

\begingroup
\setlength{\abovedisplayskip}{5pt plus 1pt minus 1pt}
\setlength{\belowdisplayskip}{5pt plus 1pt minus 1pt}
\setlength{\intextsep}{5pt}
\captionsetup{skip=6pt}
\rulesetup

\begin{figure}[H]
\plateheader{Runtime term typing}
 {\semtyp{C}{\Omega;\sigma}{t}{T}}
\noindent
\begin{minipage}[t]{0.47\linewidth}
\infrule[\ruledef{te-path}]
  {\sigma\vdash p\resolve\loc x
   \\
   \Omega;\sigma\vdash
     \sing p\capt\set p\tcoerce T
   \\
   \Omega;\sigma\vdash\set p\semleq C}
  {\semtyp{C}{\Omega;\sigma}{p}{T}}

\infrule[\ruledef{te-value}]
  {\Omega;\sigma\vdash_{\mathsf v}v:T@Q
   \\
   \Omega;\sigma\vdash\set{}\semleq C}
  {\semtyp{C}{\Omega;\sigma}{v}{T}}
\end{minipage}\hfill
\begin{minipage}[t]{0.51\linewidth}
\infrule[\ruledef{te-app}]
  {\semtyp{C_p}{\Omega;\sigma}{p}{((x:T_0)\to U)\capt C_f}
   \\
   \semtyp{C_q}{\Omega;\sigma}{q}{T_0}
   \\
   \Omega;\sigma\vdash\subst{U}{q}{x}\tcoerce T
   \andalso
   \Omega;\sigma\vdash C_p\cup C_q\semleq C}
  {\semtyp{C}{\Omega;\sigma}{p\,q}{T}}

\infrule[\ruledef{te-let}]
  {\semtyp{C}{\Omega;\sigma}{t}{T_0}
   \\
   \mathsf{Body}_{\Omega;\sigma}
      (\typ{C}{x:T_0}{u}{U})
   \\
   \Omega;\sigma\vdash U\tcoerce T
   \andalso
   \Omega;\sigma\vdash C\semleq D}
  {\semtyp{D}{\Omega;\sigma}
    {\LET x=t\ \IN u}{T}}
\end{minipage}

\caption{The complete runtime term-typing judgment.}
\Description{All rules for runtime terms with result types and use sets.}
\label{fig:runtime-term-typing}
\end{figure}

\begin{figure}[H]
\plateheader{Continuation typing}
 {\Omega;\sigma\vdash k:(U,D)\Rightarrow(T,C)}
\noindent
\begin{minipage}[t]{0.39\linewidth}
\infax[\ruledef{ke-hole}]
  {\Omega;\sigma\vdash[\,]:(T,C)\Rightarrow(T,C)}
\end{minipage}\hfill
\begin{minipage}[t]{0.58\linewidth}
\infrule[\ruledef{ke-cons}]
  {\Omega;\sigma\vdash k:(T_1,C_1)\Rightarrow(T,C)
   \\
   \mathsf{Body}_{\Omega;\sigma}
      (\typ{C_2}{x:T_0}{u}{T_2})
   \\
   \Omega;\sigma\vdash T_2\tcoerce T_1
   \andalso
   \Omega;\sigma\vdash C_0\semleq C
   \andalso
   \Omega;\sigma\vdash C_2\semleq C_1}
  {\Omega;\sigma\vdash
   (u::k):(T_0,C_0)\Rightarrow(T,C)}
\end{minipage}

\plateheader{State typing}
 {\Omega\models\langle\sigma,k,t\rangle:(T,C)}
\infrule[\ruledef{se-state}]
  {\Omega\valid\sigma
   \andalso
   \Omega;\sigma\vdash k:(U,D)\Rightarrow(T,C)
   \\
   \semtyp{D}{\Omega;\sigma}{t}{U}}
  {\Omega\models\langle\sigma,k,t\rangle:(T,C)}

\caption{The complete continuation- and state-typing judgments.}
\Description{All rules for continuation and machine-state typing.}
\label{fig:machine-typing}
\end{figure}

\endgroup

Value typing separates captured capabilities from immediate uses. A value has
an empty use set in \ruleref{te-value}, while its assigned capture set remains in the
premise. Runtime typing of an application records both operand use sets and
their common upper bound. Continuation typing carries that set through
suspended let bodies, so state typing gives one use set for the remainder of
the execution.

\begin{lemma}[Source typing interpretation]
\label{lem:source-interpretation}
If \(\typ{C}{\Gamma}{t}{T}\),
\(\rho\models^\Omega_\sigma\Gamma\), and \(\Omega\valid\sigma\), then
\[
 \semtyp{\rho(C)}{\Omega;\sigma}{\rho(t)}{\rho(T)}.
\]
\end{lemma}

The proof follows source typing.  Abstraction and let retain their body
derivations, including the use set.  Pair introductions use the
introduction types above. Applications interpret both paths. Subsumption
compiles type subtyping and subcapturing and composes the resulting coercion
and store-indexed subcapturing derivations at the outer constructor.

\section{Soundness and operational bounds}
\label{sec:soundness}

Write \(s\steps s'\) for the reflexive-transitive closure of machine steps,
and let \(\mathsf{progress}(s)\) mean that \(s\) is final or takes a step.

\begin{theorem}[Progress]
\label{thm:semantic-progress}
If \(\Omega\models s:(T,C)\), then \(\mathsf{progress}(s)\).
\end{theorem}

The runtime term-typing rule determines the proof. Paths resolve, values
return or allocate, and lets push a frame.  In the application case,
coercion action exposes a stored abstraction at the
operator location and a well-typed argument at the operand location.

Allocation changes the location scope.  We write \(T\extends U\) and
\(C\extends D\) when the right-hand object is obtained from the left by zero
or more uniform weakenings into freshly extended stores.

\begin{theorem}[One-step preservation]
\label{thm:preservation}
If \(\Omega\models s:(T,C)\) and \(s\sstep s'\), then there are a valid
annotated target store \(\Omega'\), a type \(U\), and a capture set \(D\)
such that
\[
 T\extends U,
 \qquad C\extends D,
 \qquad \Omega'\models s':(U,D).
\]
\end{theorem}

Nonallocating transitions retain \(\Omega,T,C\). Allocation obtains
\(v:T_0@Q\) from the current value typing, extends the store annotation by
\(v@Q\), and applies \ruleref{store-allocate}. All earlier typing derivations are
weakened once, including the final type and use set, after which the waiting
body is interpreted at the fresh location.

The path case uses co-resolution to convert the source path to the location
returned by \ruleref{e-path}. The let transitions move a suspended body
typing derivation into
the continuation or instantiate it with an existing location.  The
application case is the only transition in which a stored value's
assigned capture set contributes to the use set; it is detailed next.

\subsection{Beta reduction and use coverage}

Suppose an application \(p\,q\) resolves its operator to a stored closure
with assigned capture set \(Q\), and resolves its argument to location
\(y\). The suspended body typing retained by \ruleref{sem-loc-fun} has use set
\(Q\cup\set x\). Opening the body at \(y\) therefore yields use set
\(Q\cup\set y\).

Runtime term typing supplies
\[
 \Omega;\sigma\vdash\set p\semleq C_p,
 \qquad
 \Omega;\sigma\vdash\set q\semleq C_q,
 \qquad
 \Omega;\sigma\vdash C_p\cup C_q\semleq C.
\]
Since \(p\) names the closure, \ruleref{cap-fold} gives
\(Q\semleq\set p\).  Since \(q\) and \(y\) resolve to the same location,
\ruleref{cap-alias} gives \(\set y\semleq\set q\).  Hence
\[
\begin{array}{rcl}
 Q &\semleq& \set p\semleq C_p\semleq C,
 \\
 \set y &\semleq& \set q\semleq C_q\semleq C,
\end{array}
\]
and \ruleref{cap-union-elim} derives
\(Q\cup\set y\semleq C\). The reduced body consequently has the same
use set as the application. The deferred codomain coercion is
instantiated with \(y\), and runtime conversion relates the codomain opened
with \(y\) to the source result opened with \(q\).

\begin{theorem}[Finite preservation]
\label{thm:finite-preservation}
If \(\Omega\models s:(T,C)\) and \(s\steps s'\), then there are
\(\Omega',U,D\) such that \(\Omega'\) is valid for the target store,
\(T\extends U\), \(C\extends D\), and
\(\Omega'\models s':(U,D)\).
\end{theorem}

The empty annotated store validates the empty store.  Source interpretation
with the empty valuation and an identity empty continuation gives initial
state typing.

\begin{theorem}[Closed progress]
\label{thm:closed-progress}
If \(\typ{C}{\ctxempty}{t}{T}\), then
\(\mathsf{progress}(\mathsf{initial}(t))\).
\end{theorem}

\begin{theorem}[Closed finite preservation]
\label{thm:closed-preservation}
If \(\typ{C}{\ctxempty}{t}{T}\) and
\(\mathsf{initial}(t)\steps s\), then there are
\(\Omega,U,D\) such that \(\Omega\) validates the store in \(s\),
\(T\extends U\), \(C\extends D\), and
\(\Omega\models s:(U,D)\).
\end{theorem}

\begin{theorem}[Type safety]
\label{thm:type-safety}
If \(\typ{C}{\ctxempty}{t}{T}\) and
\(\mathsf{initial}(t)\steps s\), then \(\mathsf{progress}(s)\).
\end{theorem}

This is the ordinary non-stuckness result: every finite execution prefix of a
closed, well-typed program ends in a final state or a state that can step.

\subsection{Application coverage}

An application event records the two paths inspected by an application step:
\begingroup
\rulesetup
\plateheader{Application event}
 {\mathsf{application}(s\sstep s';p,q)}
\infrule[\ruledef{app-event}]
  {\sigma\vdash p\resolve\loc z
   \andalso
   \sigma\vdash q\resolve\loc y
   \andalso
   \sigma(z)=\lambda(x:T_0).t}
  {\mathsf{application}(
    \langle\sigma,k,p\,q\rangle\sstep
    \langle\sigma,k,\subst{t}{y}{x}\rangle;p,q)}
\endgroup

\begin{theorem}[Application coverage]
\label{thm:application-coverage}
Suppose \(\typ{C}{\ctxempty}{t}{T}\),
\(\mathsf{initial}(t)\steps s\), and a step from \(s\) is an application
event with operator path \(p\) and argument path \(q\). Let \(\sigma\) be the
store of \(s\). Then there are an annotation \(\Omega\) and a capture set
\(D\) such that
\[
 \Omega\valid\sigma,
 \qquad
 C\extends D,
 \qquad
 \Omega;\sigma\vdash\set p\semleq D,
 \qquad
 \Omega;\sigma\vdash\set q\semleq D.
\]
\end{theorem}

The result follows by finite preservation and inversion of
\ruleref{te-app}; continuation typing lifts the term-local coverage to the
use set \(D\). It concerns precisely the operand paths inspected by
the CK application transition.

\subsection{Capture-set bound for returned values}

The final-state relation \(s\returns v\) exposes the returned value whether
it appears directly or through a final location:
\begingroup
\rulesetup
\infrule[\ruledef{returns-value}]
  {v\ \mathsf{value}}
  {\langle\sigma,[\,],v\rangle\returns v}

\infrule[\ruledef{returns-location}]
  {\sigma(x)=v}
  {\langle\sigma,[\,],x\rangle\returns v}
\endgroup

\begin{theorem}[Returned-value capture-set bound]
\label{thm:returned-capture-bound}
Suppose \(\typ{C}{\ctxempty}{t}{T}\),
\(\mathsf{initial}(t)\steps s\), and \(s\) is final. Let \(\sigma\) be the
store of \(s\). Then there are an annotation \(\Omega\), a shape type \(R\),
capture sets \(C',D,Q\), a returned value \(v\), and a type \(U\) such that \(Q\) is the assigned
capture set of \(v\) and
\[
 \begin{gathered}
 \Omega\valid\sigma,
 \qquad T\extends R\capt C',
 \qquad C\extends D,
 \\
 s\returns v,
 \qquad
 \Omega;\sigma\vdash_{\mathsf v}v:U@Q,
 \qquad
 \Omega;\sigma\vdash Q\semleq C'.
 \end{gathered}
\]
\end{theorem}

For a direct value, inversion of the empty continuation identifies the
runtime value's indices with the final type and capture set, after which the
conclusion follows from value typing. For a returned location, annotated
store typing supplies the value and its assigned capture set;
\ruleref{cap-fold} and the coercion retained by path typing relate that set
to \(C'\).

Type safety, application coverage, and the returned-value capture-set bound have
distinct conclusions. Type safety rules out stuck states. Application
coverage bounds the paths inspected by application. The final result bounds
the assigned capture set of a returned value. The calculus has no primitive capability
operation or effect event, so these results do not constitute a theorem about
I/O or another external effect semantics.

\section{Limitations and next steps}

The calculus has immutable stores, monadic normal form, stable paths,
dependent functions, singleton types, abstract type and capture-set members, and
dependent pairs.  It omits intersections, recursive self types, mutable
state, higher-kinded types, and a module initialization discipline.  These
omissions keep the proof independent of the narrowing and inert-type
infrastructure used by larger DOT calculi.

The dynamic observations in this development are application operands and
returned values.  Adding primitive capabilities would require syntax and an
event semantics for their operations, together with a theorem connecting
those events to the use set. The metatheoretic coercions would remain proof
objects.

A next extension can add intersections in a controlled fragment and
determine which distributivity and member-merging principles are needed by
Scala encodings.  Recursive objects require a precise account of
initialization and stable self paths.  The general pair proof identifies the
main semantic constraint: coercion action relies on members referring only to
older store locations.  Mutable state or recursive allocation would break
this stratification and require another model, such as a step-indexed
relation.

The mechanization proves bound consistency and its preservation by member
subtyping.  Full context formation and typing regularity remain open because
they require the opening and path-substitution metatheory needed to recover
well-formed synthesized signatures from arbitrary typing derivations.

\section{Related work}

The calculus follows the path-dependent tradition initiated by foundations
for Scala and DOT \citep{amin2014foundations,rompf2016dot}.  Later DOT and
pDOT proofs develop precise typing, inert contexts, and path replacement for
substantially richer calculi
\citep{rapoport2017simple,rapoport2019pdot}.  The restricted type language
here permits a direct store-indexed interpretation.  Step-indexed logical
relations provide another route to full DOT and mutable or recursive
features \citep{giarrusso2020gdot}; the append-only allocation order suffices
for the present system.

Interpreting subtyping as explicit evidence has a long history
\citep{breazutannen1991coercion,crary2000inclusive}.  System FC makes equality
evidence part of an explicitly typed intermediate language
\citep{sulzmann2007systemfc}.  Our coercions remain semantic objects used in
the metatheory, while the source language contains neither coercion terms nor
modalities.  An explicit intermediate elaboration may become useful when
connecting the calculus to a realistic compiler
\citep{martres2023compilation}.

Capturing types and subcapturing follow Capturing Types
\citep{boruch2023capture}. The use-set typing judgment follows System Capless,
which provides explicit universal and existential capture quantification.
System Cappy gives a lightweight surface syntax and a type-preserving
translation to Capless, avoiding explicit capture variables in ordinary
programs \citep{xu2025box}. Abstract capture-set members
provide path-dependent capture-set abstraction here. Application coverage
relates use sets to the paths inspected during evaluation, while the
returned-value theorem bounds the assigned capture set of the returned value.

Type and capture-set members both use intervals. Their variance follows the
standard order on lower and upper bounds; consistency premises exclude
inverted intervals at formation and selection sites.  More expressive
higher-order interval systems raise additional questions outside the current
calculus \citep{stucki2021intervals}.

\section{Conclusion}

\lambdapcc{} combines path-dependent types and capture checking through one
dependent-pair mechanism. Abstract type and capture-set members share path
lookup and interval subtyping; term paths additionally support root
contraction.  Store-indexed coercions account for explicit transitivity,
path selections, dependent functions, and the general pair rule. The runtime
typing invariant preserves both result types and use sets while retaining the
assigned capture set of each stored value. The mechanized development proves
progress, preservation, and non-stuckness, covers both paths inspected by
every reachable application, and bounds the assigned capture set of every
returned value by the capture set of its result type.

\bibliographystyle{ACM-Reference-Format}
\bibliography{../../../references}

\end{document}
